% !TeX root = ../main.tex
% !TeX encoding = UTF-8
% !TeX spellcheck = none

\section{Problem Definition}\label{sec:problem-definition}

We partition the study area into \(N\) regular grid cells and treat each cell as a forecasting node. At time \(t\), the regional flow matrix is
\begin{equation}
\label{eq:flow}
X_t = \bigl[x_t^{\mathrm{in}},\,x_t^{\mathrm{out}}\bigr]
\in \mathbb{R}^{N\times 2},
\end{equation}

where \(x_t^{\mathrm{in}},x_t^{\mathrm{out}}\in\mathbb{R}^{N}\) are the inflow and outflow during the sampling interval. They count trips ending and starting in each region, respectively~\cite{ref32}. A single model forecasts both channels, ordered as inflow followed by outflow. Sequence tensors use time, node, and channel order, with the batch dimension omitted.

Given a historical observation sequence ending at time \(t\) with length \(T\),
\begin{equation}
\label{eq:history}
X_{t-T+1:t} = \bigl[X_{t-T+1},\ldots,X_t\bigr]
\in \mathbb{R}^{T\times N\times 2},
\end{equation}

the goal is to jointly forecast inflow and outflow over the next \(H\) time steps:
\begin{equation}
\label{eq:future}
Y_{t+1:t+H} = \bigl[X_{t+1},\ldots,X_{t+H}\bigr]
\in \mathbb{R}^{H\times N\times 2}.
\end{equation}

Here, \(T\) and \(H\) are the history length and number of forecast steps. With a 5-minute sampling interval, we set \(T=H=12\), using the past 60 minutes to forecast the next 60 minutes.

We formulate regional traffic flow forecasting as learning the spatiotemporal mapping
\begin{equation}
\label{eq:mapping}
\widehat{Y}_{t+1:t+H}
= f_{\theta}\!\left(X_{t-T+1:t},C_t;\mathcal{S},\mathcal{B}\right)
\in \mathbb{R}^{H\times N\times 2},
\end{equation}

Here, \(\theta\) denotes the learnable parameters. The temporal information \(C_t\in\mathbb{R}^{T\times 2}\) contains time of day and day of week for the historical window and is shared across nodes. The fixed spatial supports \(\mathcal{S}\) include self-relations, flow-specific directed time-lagged associations, and candidate cross-region neighborhoods. The lookup table \(\mathcal{B}\) stores the routine-flow reference estimated from training data. Forecasts use observations up to \(t\); \(\widehat{Y}_{t+1:t+H}\) is expressed on the original flow scale.

For each historical time \(u\in\{t-T+1,\ldots,t\}\), we retrieve the corresponding routine-flow reference from \(\mathcal{B}\):
\begin{equation}
\label{eq:baseline}
b_u=\bigl[b_u^{\mathrm{in}},\,b_u^{\mathrm{out}}\bigr]
\in\mathbb{R}^{N\times 2}.
\end{equation}

We estimate the reference from training observations grouped by node, channel, weekday or weekend, and time-of-day slot. The flow deviation is
\begin{equation}
\label{eq:deviation}
D_u=X_u-b_u\in\mathbb{R}^{N\times 2}.
\end{equation}

The deviation \(D_u\) measures the difference between observed flow and its routine level. Section~\ref{sec:methodology} uses this reference with historical trajectories to represent trend deviation and short-term fluctuation.

The model corrects the base prediction when traffic deviates from routine patterns. It combines recent node dynamics with cross-region time-lagged associations, then selects historical messages using source and target backgrounds and the forecast horizon. Section~\ref{sec:methodology} gives the residual refinement mechanism.
